Choose \(\mathbf M^0\in\mathcal C_{\mathrm{sat},D}(\mathcal P)\).  District factorization gives the finite nonfocal product law \(\nu_{-D}^{0,s}\).  Under \(\mathsf{ASep}_D(\mathcal P)\), every focal slice lies in the linear space \(R_D(m_D)\), so
\[
\bar\chi_C^{s,\mathbf M^0}
=\sum_{\omega_{-D}}
\nu_{-D}^{0,s}(\omega_{-D})
\chi_{C,\omega_{-D}}
\in R_D(m_D).
\tag{1}
\]
The same baseline works for every cell and domain, proving \(\mathsf{BASep}_D\).  Composition with the structural-separation lemma gives the full hierarchy.

For strictness, take independent fair binary roots \(X,Z\) feeding the deterministic binary child \(Y=\mathbf 1\{X=Z\}\), and let \(D=\{X\}\).  Supply only the complete target-domain output law of \(Y\).  Then \(A_D=\mathbf 1^{\mathsf T}\) and
\[
R_D(m_D)=\operatorname{span}\{(1,1)\}.
\]
The two focal slices of \(\{Y=1\}\) are \((1,0)\) and \((0,1)\), so active-slice separation fails.  Their average under the fair nonfocal law of \(Z\) is \((1/2,1/2)\in R_D(m_D)\); the complementary output cell is analogous.  Focal saturation permits every law of \(X\) while the fair law of \(Z\) and the structural table of \(Y\) stay fixed, so this single compatible baseline proves \(\mathsf{BASep}_D\).  Thus the second implication is strict on the original binary alphabets.